\chapter{System Overview}
\label{ch:system}

This chapter presents the overall architecture of \projectshort{}. The system is organized as a three-layer preprocessing pipeline---scene analysis, adaptive mathematical denoising, and optional lightweight deep refinement---wrapped by a web-based evaluation platform. Both a command-line batch workflow and an interactive web application share the same core algorithm library.

\section{Design Goals}
\label{sec:design_goals}

The platform was designed for local audio preprocessing research in an ASR-oriented setting (\chapref{ch:introduction}). Four engineering goals guided architectural decisions:

\begin{itemize}
  \item \textbf{Interpretability.} Processing stages should remain inspectable: routing decisions log explicit algorithm chains and human-readable reasons; OM-LSA/MCRA gains and noise PSD updates can be traced frame-by-frame; residual waveforms are exported for diagnosis.
  \item \textbf{Adaptivity.} Enhancement should respond to scene characteristics---band-wise SNR imbalance, tonal structure, reverberation, and local SNR variability---via both internal bin-level adaptation and external algorithm routing.
  \item \textbf{Efficiency.} Deployment should not require large-model inference. Mathematical denoising carries the main computational load; optional refinement uses a TinyResidualRefiner with only thousands of parameters. DeepFilterNet3 is confined to offline teacher generation during distillation training.
  \item \textbf{Evaluability.} Every processed utterance should yield reproducible objective metrics (SNR, residual RMS, kurtosis), 24+ diagnostic plots, synchronized A/B audio comparison, and ABX blind listening statistics.
\end{itemize}

These goals motivate a hybrid design rather than a monolithic end-to-end neural enhancer or a single fixed classical denoiser.

\section{Three-Layer Architecture}
\label{sec:architecture}

\figref{fig:system_architecture} summarizes the end-to-end pipeline. Each layer addresses a distinct sub-problem while passing structured artifacts to the next stage.


\paragraph{Layer 1: Full-dimensional scene analysis.}
For each benchmark scenario, clean speech, isolated noise, and the noisy mixture are analyzed jointly. Extracted statistics include time-domain RMS and zero-crossing rate (ZCR), Welch power spectral density, band-wise SNR (low/mid/high), local SNR heatmaps, and residual descriptors (kurtosis, spectral flatness). These quantities feed the ANASS parameter guidance framework (\chapref{ch:scene}) and inform routing thresholds (\chapref{ch:routing}).

\paragraph{Layer 2: Adaptive mathematical denoising.}
The rule router reads precomputed scene metrics (or accepts manual overrides) and selects one or more algorithms from a library of seven modules. Chains are composable: for example, subspace projection may precede chirp notch filtering, or adaptive notch filtering may precede NMF suppression. All chains terminate in a shared OM-LSA/MCRA back-end (\chapref{ch:algorithms}).

\paragraph{Layer 3: Lightweight deep refinement.}
When enabled, the TinyResidualRefiner (\chapref{ch:dl}) takes the routed STFT magnitude and the original noisy STFT magnitude as input and predicts a multiplicative residual mask. This stage closes part of the quality gap to the DeepFilterNet3 teacher without invoking the teacher at inference.

\paragraph{Evaluation and web platform.}
Enhanced audio, residual signals, JSON metrics, and Plotly plot specifications are exposed through the web application (\chapref{ch:web}) for interactive review and ABX testing.

\begin{figure}[H]
  \centering
  \resizebox{0.98\textwidth}{!}{%
    % Three-layer system architecture (included from Chapter 3)
% Inter-layer spacing anchored to fit-box bottoms (not inner nodes)
\begin{tikzpicture}[
  node distance=0.85cm and 0.75cm,
  block/.style={
    draw=shublue,
    line width=0.6pt,
    fill=white,
    rounded corners=2pt,
    minimum height=0.9cm,
    minimum width=2.5cm,
    align=center,
    font=\footnotesize,
    inner sep=4pt,
  },
  wideblock/.style={
    block,
    minimum width=11.5cm,
  },
  smallblock/.style={
    block,
    minimum width=2.2cm,
    minimum height=0.8cm,
    font=\scriptsize,
  },
  layerbox/.style={
    draw=shublue!70,
    line width=0.5pt,
    rounded corners=6pt,
    inner xsep=0.55cm,
    inner ysep=0.65cm,
  },
  layertitle/.style={
    font=\bfseries\small,
    text=shublue,
    align=center,
  },
  arr/.style={
    -{Stealth[length=2.2mm, width=1.8mm]},
    line width=0.55pt,
    draw=shublue,
  },
  dashedarr/.style={
    arr,
    dashed,
  },
  note/.style={
    font=\scriptsize,
    text=shugray,
    align=center,
  },
]

% --- vertical gaps between layer boxes ---
\pgfmathsetlengthmacro{\interlayer}{2.85cm}   % box bottom to next layer top row
\pgfmathsetlengthmacro{\titleabove}{0.38cm}   % title above each fit box

% =============================================================================
% Layer 1: Scene Analysis
% =============================================================================
\node[block] (clean) {Clean\\Reference};
\node[block, right=of clean] (noise) {Noise\\Component};
\node[block, right=of noise] (mix) {Noisy\\Mixture};

\node[wideblock, below=0.95cm of noise] (metrics) {%
  Scene Metrics Extraction\\
  \scriptsize (RMS, ZCR, band SNR, local SNR, flatness)};

\node[wideblock, below=of metrics] (anass) {%
  ANASS Parameter Guidance \& Full-Dimensional Reports};

\begin{scope}[on background layer]
  \node[layerbox, fill=shulayer1, fit=(clean)(mix)(anass)] (layer1box) {};
\end{scope}

\node[layertitle, above=\titleabove of layer1box.north, anchor=south] (layer1title)
  {Layer 1: Full-Dimensional Scene Analysis};

\draw[arr] (clean.south) -- ++(0,-0.28) -| ([xshift=-0.15cm]metrics.north);
\draw[arr] (noise.south) -- (metrics.north);
\draw[arr] (mix.south) -- ++(0,-0.28) -| ([xshift=0.15cm]metrics.north);
\draw[arr] (metrics) -- (anass);

% =============================================================================
% Layer 2: Adaptive Math Denoising  (anchored below Layer 1 box)
% =============================================================================
\coordinate (layer2anchor) at ([yshift=-\interlayer]layer1box.south);

\node[block, anchor=north, minimum width=2.8cm] (router) at (layer2anchor) {Rule Router};
\node[block, right=of router, minimum width=3.6cm] (chain) {%
  Algorithm Chain\\[1pt]\scriptsize (7 modules, composable)};
\node[block, right=of chain, minimum width=2.8cm] (omlsa) {%
  OM-LSA/MCRA\\Back-end};

\begin{scope}[on background layer]
  \node[layerbox, fill=shulayer2, fit=(router)(omlsa)] (layer2box) {};
\end{scope}

\node[layertitle, above=\titleabove of layer2box.north, anchor=south] (layer2title)
  {Layer 2: Adaptive Mathematical Denoising};

\draw[arr] (anass.south) -- ++(0,-0.55) -| (router.north);
\draw[arr] (router) -- (chain);
\draw[arr] (chain) -- (omlsa);

% =============================================================================
% Layer 3: Lightweight Deep Refinement  (anchored below Layer 2 box)
% =============================================================================
\coordinate (layer3anchor) at ([yshift=-\interlayer]layer2box.south);

\node[smallblock, anchor=north] (mathout) at (layer3anchor) {Routed\\Output};
\node[block, right=of mathout, minimum width=3.8cm] (student) {%
  TinyResidualRefiner\\[1pt]\scriptsize (optional distill refine)};
\node[smallblock, right=of student] (refined) {Refined\\Output};

\node[note, below=0.12cm of student.south] (studentnote)
  {inputs: routed STFT + noisy STFT};

\begin{scope}[on background layer]
  \node[layerbox, fill=shulayer3, fit=(mathout)(refined)(studentnote)] (layer3box) {};
\end{scope}

\node[layertitle, above=\titleabove of layer3box.north, anchor=south] (layer3title)
  {Layer 3: Lightweight Deep Refinement};

\draw[arr] (omlsa.south) -- ++(0,-0.55) -| (mathout.north);
\draw[arr] (mathout) -- (student);
\draw[arr] (student) -- (refined);

% Offline teacher in the gap between Layer 3 and Layer 4
\node[smallblock, below=1.15cm of layer3box.south, draw=shugray, dashed,
      anchor=north] (teacher) {DeepFilterNet3\\(offline teacher)};

\draw[dashedarr, draw=shugray]
  (teacher.north) -- node[note, right=0.15cm] {distill train} (student.south);

% =============================================================================
% Layer 4: Evaluation & Web Platform  (anchored below teacher)
% =============================================================================
\coordinate (layer4anchor) at ([yshift=-\interlayer]teacher.south);

\node[block, anchor=north, minimum width=2.4cm] (metrics_out) at (layer4anchor) {Metrics\\Snapshot};
\node[block, right=of metrics_out, minimum width=2.5cm] (plots) {24+ Diagnostic\\Plots};
\node[block, right=of plots, minimum width=2.4cm] (audio) {A/B Audio\\Compare};
\node[block, right=of audio, minimum width=2.4cm] (abx) {ABX Blind\\Listening};

\begin{scope}[on background layer]
  \node[layerbox, fill=shulayer4, fit=(metrics_out)(abx)] (layer4box) {};
\end{scope}

\node[layertitle, above=\titleabove of layer4box.north, anchor=south] (layer4title)
  {Evaluation \& Web Platform};

\draw[arr] (refined.south) -- ++(0,-0.55) -| (plots.north);
\draw[arr] (refined.south) -- ++(0,-0.55) -| (audio.north);
\draw[arr] (plots) -- (metrics_out);
\draw[arr] (audio) -- (abx);

\end{tikzpicture}
%
  }
  \caption{Three-layer architecture of \projectshort{} with evaluation platform. Solid arrows denote runtime data flow; the dashed path indicates offline teacher usage during distillation training only.}
  \label{fig:system_architecture}
\end{figure}
\section{Data Flow}
\label{sec:data_flow}

\subsection{Batch Pipeline (CLI)}

The offline workflow centered on \texttt{run\_pipeline.py} follows six steps:

\begin{enumerate}
  \item \textbf{Input.} A mono WAV file from \texttt{noisy\_testset/} (e.g., \texttt{scene01\_white\_snr15.wav}).
  \item \textbf{Routing.} \texttt{router/rule\_router.py} loads \texttt{analysis/scene\_details/\{scene\}\_metrics.json}, checks \texttt{config/scene\_method\_overrides.json}, and returns a \texttt{RouteDecision} (chain + reason).
  \item \textbf{Enhancement.} Each method in the chain is applied sequentially via \texttt{algorithms/*.py}, producing \texttt{*\_routed.wav} in \texttt{outputs/routed/}.
  \item \textbf{Optional distillation.} \texttt{distill/infer\_student\_refine.py} loads \texttt{student\_runD.pt} and overwrites the routed output with a refined version in \texttt{outputs/distill\_refined/}.
  \item \textbf{Logging.} Route metadata is written to \texttt{outputs/route\_log.json}.
  \item \textbf{Evaluation.} \texttt{eval/batch\_eval\_metrics.py} and \texttt{eval/search\_best\_method\_per\_scene.py} compute RMS errors and search per-scene best methods.
\end{enumerate}

\subsection{Web Application Flow}

\figref{fig:web_dataflow} depicts the interactive workflow. The user uploads audio through the React frontend; the FastAPI backend creates a task, stores the original WAV, and queues background processing.

\begin{figure}[H]
  \centering
  \resizebox{0.92\textwidth}{!}{%
    % Web application data-flow diagram (included from Chapter 3)
% Three-column bottom row; sequence strip centered below entire processing tier
\begin{tikzpicture}[
  node distance=0.9cm and 1.65cm,
  actor/.style={
    draw=shublue,
    line width=0.55pt,
    fill=shulayer1,
    rounded corners=3pt,
    minimum width=2.2cm,
    minimum height=0.95cm,
    align=center,
    font=\footnotesize,
    inner sep=4pt,
  },
  service/.style={
    draw=shublue,
    line width=0.55pt,
    fill=white,
    rounded corners=3pt,
    minimum width=3.0cm,
    minimum height=1.05cm,
    align=center,
    font=\footnotesize,
    inner sep=4pt,
  },
  store/.style={
    draw=shublue,
    line width=0.55pt,
    fill=shulayer2,
    rounded corners=3pt,
    minimum height=0.95cm,
    minimum width=2.6cm,
    align=center,
    font=\footnotesize,
    inner sep=4pt,
  },
  seqbox/.style={
    draw=shublue!60,
    line width=0.45pt,
    fill=shulayer1,
    rounded corners=4pt,
    inner sep=6pt,
    font=\scriptsize,
    text=shugray,
    align=center,
  },
  arr/.style={
    -{Stealth[length=2.2mm, width=1.7mm]},
    line width=0.55pt,
    draw=shublue,
  },
  lbl/.style={
    font=\scriptsize,
    text=shugray,
    fill=white,
    inner sep=2pt,
    align=center,
  },
]

\pgfmathsetlengthmacro{\rowgap}{2.15cm}
\pgfmathsetlengthmacro{\loopheight}{1.15cm}
\pgfmathsetlengthmacro{\seqgap}{1.75cm}

% =============================================================================
% Top row (client tier)
% =============================================================================
\node[actor] (user) {User /\\Browser};
\node[service, right=of user] (frontend) {React Frontend\\Upload, Charts, ABX};
\node[service, right=of frontend] (api) {FastAPI Backend\\REST + Tasks};

\draw[arr]
  (frontend.north) -- ++(0,\loopheight)
  -| node[lbl, above, pos=0.5, yshift=-2pt] {poll status} (api.north);

\draw[arr] (user) -- node[lbl, above] {HTTPS} (frontend);
\draw[arr] (frontend) -- node[lbl, above] {/api/*} (api);

% =============================================================================
% Bottom row (processing tier)
% =============================================================================
\node[store, below=\rowgap of frontend, anchor=north] (files) {%
  WAV / JSON\\Storage};
\node[service, below=\rowgap of api, anchor=north] (runner) {%
  Denoise Runner\\+ Distill Refiner};
\node[service, right=of runner] (algo) {%
  Algorithm Library\\+ Rule Router};

\draw[arr] (api.south) -- node[lbl, right, pos=0.42] {queue} (runner.north);
\draw[arr] (runner.east) -- node[lbl, above] {invoke} (algo.west);

% read/write: route below the boxes (avoids crossing Algorithm Library)
\draw[arr]
  (runner.south) -- ++(0,-0.55)
  -| node[lbl, below, pos=0.25, yshift=2pt] {read / write} (files.south);

% persist: drop below API, then enter Storage from above
\draw[arr]
  (api.south) -- ++(0,-0.5)
  -| node[lbl, above, pos=0.3] {persist} (files.north);

% =============================================================================
% Sequence strip — centered below the full bottom row
% =============================================================================
\coordinate (bottomrow) at ($(files.south)!0.5!(algo.south)$);

\node[seqbox,
      below=\seqgap of bottomrow,
      anchor=north,
      minimum width=13.0cm] (seqstrip) {%
  \textbf{\textcolor{shublue}{Processing sequence:}}\quad
  \textbf{1.}~Upload
  $\rightarrow$ \textbf{2.}~Process
  $\rightarrow$ \textbf{3.}~Denoise
  $\rightarrow$ \textbf{4.}~Optional refine
  $\rightarrow$ \textbf{5.}~Metrics / plots
  $\rightarrow$ \textbf{6.}~ABX};

\end{tikzpicture}
%
  }
  \caption{Web application data flow. The frontend polls task status while the backend orchestrates denoising, optional refinement, and artifact generation.}
  \label{fig:web_dataflow}
\end{figure}

Upon \texttt{POST /api/process}, the backend invokes \texttt{denoise\_runner.run\_denoise()} with the selected method (\texttt{auto}, explicit algorithm name, \texttt{noisereduce}, \texttt{wiener}, or \texttt{deepfilter}). If \texttt{run\_distill\_refine=true}, \texttt{distill\_refiner.refine\_with\_student()} post-processes the denoised WAV. Finally, \texttt{metrics\_service.build\_metrics\_and\_plots()} generates SNR statistics, residual descriptors, and 24+ grouped Plotly figures stored as JSON for the chart center.

\subsection{Artifact Types}

\begin{table}[H]
  \centering
  \caption{Primary runtime artifacts produced by the system.}
  \label{tab:artifacts}
  \begin{tabular}{lll}
    \toprule
    \textbf{Artifact} & \textbf{Location / Format} & \textbf{Purpose} \\
    \midrule
    Original WAV       & \texttt{webapp/data/uploads/}     & Input preservation \\
    Denoised WAV       & \texttt{webapp/data/processed/}   & Enhanced output \\
    Residual WAV       & \texttt{webapp/data/processed/}   & $x - \hat{x}$ diagnosis \\
    Metrics JSON       & \texttt{*\_metrics.json}          & SNR, RMS, kurtosis \\
    Plots JSON         & \texttt{*\_plots.json}            & Plotly chart specs \\
    Route log          & \texttt{route\_log.json}          & Chain + reason audit \\
    Student checkpoint & \texttt{distill/checkpoints/}     & Refiner inference \\
    \bottomrule
  \end{tabular}
\end{table}

\section{Repository Structure}
\label{sec:repo_structure}

Table~\ref{tab:repo_layout} maps major directories to functional roles. Separation between \texttt{analysis/}, \texttt{analysis/denoise\_selection/}, and \texttt{noisy\_testset/} keeps research artifacts, executable code, and benchmark audio distinct.

\begin{table}[H]
  \centering
  \caption{Repository layout and responsibilities.}
  \label{tab:repo_layout}
  \begin{tabularx}{\textwidth}{l X}
    \toprule
    \textbf{Path} & \textbf{Description} \\
    \midrule
    \texttt{noisy\_testset/} & 15 noisy scenes + \texttt{clean\_reference.wav} \\
    \texttt{analysis/reports/} & Full-dimensional scene analysis reports (Markdown) \\
    \texttt{analysis/scene\_details/} & Per-scene JSON metrics consumed by the router \\
    \texttt{analysis/priority\_plots/} & Key diagnostic figures referenced in reports \\
    \texttt{analysis/denoise\_selection/algorithms/} & Seven denoising modules + shared STFT utilities \\
    \texttt{analysis/denoise\_selection/router/} & Scene-aware rule router \\
    \texttt{analysis/denoise\_selection/distill/} & Pair building, training, inference, checkpoints \\
    \texttt{analysis/denoise\_selection/webapp/} & FastAPI backend + React frontend \\
    \texttt{analysis/denoise\_selection/eval/} & Batch metrics and best-method search \\
    \texttt{report/} & This LaTeX project report \\
    \bottomrule
  \end{tabularx}
\end{table}

\section{Software Stack}
\label{sec:software_stack}

Table~\ref{tab:software_stack} lists the primary technologies. Python 3.11+ with NumPy/SciPy handles signal processing; PyTorch supports distillation only; the web tier uses mainstream async API and SPA frameworks for reproducible deployment.

\begin{table}[H]
  \centering
  \caption{Main software components.}
  \label{tab:software_stack}
  \begin{tabular}{lll}
    \toprule
    \textbf{Layer} & \textbf{Technology} & \textbf{Role} \\
    \midrule
    Scene analysis   & Python, NumPy, SciPy, Matplotlib & Metrics, reports, plots \\
    Denoising core   & Python modules (\texttt{algorithms/}) & STFT, OM-LSA/MCRA, specialized modules \\
    Routing          & Python (\texttt{router/})           & Scene-aware chain selection \\
    Distillation     & PyTorch                             & TinyResidualRefiner train/infer \\
    Teacher (offline)& DeepFilterNet3 CLI                  & Teacher WAV generation \\
    Backend API      & FastAPI, Uvicorn, python-multipart & REST endpoints, background tasks \\
    Frontend UI      & React, TypeScript, Vite, Plotly.js  & Upload, charts, ABX, history \\
    Audio I/O        & \texttt{soundfile}, optional \texttt{ffmpeg} & WAV/MP3/FLAC support \\
    Containerization & Docker Compose                      & One-command local deployment \\
    \bottomrule
  \end{tabular}
\end{table}

\section{Operating Modes}
\label{sec:operating_modes}

The same algorithm library supports three operating modes:

\begin{enumerate}
  \item \textbf{Automatic routing (\texttt{method=auto}).} Recommended default. The router selects a chain from scene metrics; suitable for benchmark reproduction and unknown user uploads whose filenames match benchmark conventions.
  \item \textbf{Explicit algorithm selection.} Users force a single module or baseline (\texttt{omlsa}, \texttt{noisereduce}, \texttt{wiener}, \texttt{deepfilter}) for ablation studies.
  \item \textbf{Auto + distill refinement.} Automatic routing followed by TinyResidualRefiner; balances quality and edge efficiency when the checkpoint is available.
\end{enumerate}

If DeepFilterNet CLI is unavailable, explicit \texttt{deepfilter} mode falls back to OM-LSA/MCRA; distillation inference silently skips when the checkpoint file is missing.

\section{Benchmark Scenarios}
\label{sec:benchmark}

Evaluation uses 15 single-channel scenes paired with a shared clean reference utterance. Table~\ref{tab:benchmark_scenes} summarizes noise types and nominal SNR levels. Together they span stationary colored noise, tonal hum, impulsive events, speech-like interference, and reverberation---conditions representative of mobile ASR front-end challenges.

\begin{table}[H]
  \centering
  \caption{Benchmark scenarios in \texttt{noisy\_testset/}.}
  \label{tab:benchmark_scenes}
  \small
  \begin{tabular}{clll}
    \toprule
    \textbf{ID} & \textbf{Filename} & \textbf{Noise type} & \textbf{SNR (dB)} \\
    \midrule
    01 & \texttt{scene01\_white\_snr15.wav}              & White noise        & 15 \\
    02 & \texttt{scene02\_white\_snr5.wav}               & White noise        & 5  \\
    03 & \texttt{scene03\_pink\_snr10.wav}               & Pink noise         & 10 \\
    04 & \texttt{scene04\_brown\_snr8.wav}               & Brown noise        & 8  \\
    05 & \texttt{scene05\_hum50\_snr12.wav}              & 50\,Hz power hum   & 12 \\
    06 & \texttt{scene06\_hum60\_snr12.wav}              & 60\,Hz power hum   & 12 \\
    07 & \texttt{scene07\_highfreq\_hiss\_snr10.wav}     & High-frequency hiss& 10 \\
    08 & \texttt{scene08\_lowfreq\_rumble\_snr10.wav}     & Low-frequency rumble& 10 \\
    09 & \texttt{scene09\_chirp\_interference\_snr8.wav} & Chirp interference & 8  \\
    10 & \texttt{scene10\_impulsive\_clicks\_snr9.wav}   & Impulsive clicks   & 9  \\
    11 & \texttt{scene11\_intermittent\_bursts\_snr7.wav}& Intermittent bursts& 7  \\
    12 & \texttt{scene12\_babble\_like\_snr5.wav}        & Babble-like noise  & 5  \\
    13 & \texttt{scene13\_street\_combo\_snr6.wav}       & Street mixture     & 6  \\
    14 & \texttt{scene14\_office\_fan\_snr9.wav}          & Office fan noise   & 9  \\
    15 & \texttt{scene15\_reverb\_plus\_hiss\_snr10.wav} & Reverberation + hiss& 10 \\
    \bottomrule
  \end{tabular}
\end{table}

Priority optimization scenarios identified by scene analysis---including \texttt{scene02}, \texttt{scene09}, \texttt{scene10}, \texttt{scene11}, and \texttt{scene12}---receive higher weight during distillation training and are emphasized in ablation discussions (\chapref{ch:experiments}).

\section{Chapter Summary}
\label{sec:system_summary}

\projectshort{} integrates analysis, adaptive classical enhancement, optional distilled refinement, and interactive evaluation into a cohesive system suitable for local ASR preprocessing research. \chapref{ch:scene} details how Layer~1 statistics constrain algorithm design; \chapref{ch:algorithms} and \chapref{ch:routing} expand Layer~2; \chapref{ch:dl} and \chapref{ch:web} complete Layers~3 and the evaluation shell respectively.
